% ============================================
% METHODS CHAPTER
% ============================================

\section{Methods}
\label{sec:methods}

Building on \Cref{def:gen-error} from the introduction, we now
develop the theoretical framework.

\subsection{Mathematical Framework}

\begin{theorem}[PAC Learning Bound]
\label{thm:main}
Let $\mathcal{H}$ be a hypothesis class with VC dimension $d$.
For any $\delta > 0$, with probability at least $1 - \delta$:
\begin{equation}
\label{eq:pac-bound}
R(h) \leq \hat{R}(h) + \sqrt{\frac{d \log(2n/d) + \log(4/\delta)}{n}}
\end{equation}
where $\hat{R}(h)$ is the empirical risk.
\end{theorem}

\begin{proof}
The proof follows from uniform convergence bounds. See \cite{shalev2014}
for complete details.
\end{proof}

\subsection{Matrix Operations}

Consider the covariance matrix $\Sigma \in \R^{d \times d}$:
\begin{equation}
\Sigma = \begin{pmatrix}
\sigma_1^2 & \rho_{12} & \cdots & \rho_{1d} \\
\rho_{21} & \sigma_2^2 & \cdots & \rho_{2d} \\
\vdots & \vdots & \ddots & \vdots \\
\rho_{d1} & \rho_{d2} & \cdots & \sigma_d^2
\end{pmatrix}
\end{equation}

The eigenvalue decomposition yields $\Sigma = U \Lambda U^\top$.

\subsection{Algorithm}

\begin{algorithm}[H]
\caption{Gradient Descent for Empirical Risk Minimization}
\label{alg:gd}
\KwIn{Dataset $\mathcal{D}$, learning rate $\eta$, iterations $T$}
\KwOut{Optimized parameters $\theta^*$}
Initialize $\theta_0$ randomly\;
\For{$t = 1$ \KwTo $T$}{
    Compute gradient: $g_t \gets \nabla_\theta \hat{R}(\theta_{t-1})$\;
    Update: $\theta_t \gets \theta_{t-1} - \eta \cdot g_t$\;
    \If{$\norm{g_t} < \epsilon$}{
        \textbf{break}\;
    }
}
\Return $\theta_T$\;
\end{algorithm}

See \Cref{alg:gd} for the complete procedure.

\subsection{Aligned Equations}

The loss function and its gradient:
\begin{align}
\mathcal{L}(\theta) &= \frac{1}{n} \sum_{i=1}^{n} \ell(h_\theta(x_i), y_i) + \lambda \norm{\theta}^2 \label{eq:loss} \\
\nabla \mathcal{L}(\theta) &= \frac{1}{n} \sum_{i=1}^{n} \nabla_\theta \ell(h_\theta(x_i), y_i) + 2\lambda\theta \label{eq:grad}
\end{align}

\Cref{eq:loss,eq:grad} define the optimization objective.
